\chapter{数列}
    % \section{実数を要素とする有界集合の上限と下限の存在(考察)}
    % 	$A\in\realNumbers$を上に有界な集合とし、$A$の上限$\sup A$の存在を確かめる。
    % 	同様にして下に有界な集合に対する下限の存在も確かめられる。
    % 	まず$a_0 \in A$を任意に選ぶ。
    % 	$A$は上に有界だから、次のような$a_1$が存在する。
    % 	\[ a_0 < b_0,\; \forall a \in A,\; a<b_0 \]
    % 	数列$\{a_i\},\{b_i\}$を次の手順で生成する:
    % 	$(a_i+b_i)/2$が$A$の上界であれば$a_{i+1} = a_i,\; b_{i+1} = (a_i+b_i)/2$、そうでなければ$a_{i+1} = (a_i+b_i)/2,\; b_{i+1} = b_i$とする。
    % 	$c_i = (a_i+b_i)/2$とすると数列$\{c_i\}$はCauchy列となるので収束する。
    \section{単調性}
        \subsection{\texorpdfstring{$\left(1-\frac{t}{n}\right)^n\;(t\in[0,n])$}{(1-t/n)\string^n (t∈[0,n])}は\texorpdfstring{$n$}{n}に関して単調増加}
            \label{(1-t/n)^nは単調増加}
            \begin{proof}
                \quad\par
                $a_n \coloneqq \left(1-\frac{t}{n}\right)^n$とする。
                $t=0,n$の場合は単調増加性は明らか。以下$0<t<n$とする。
                \begin{align*}
                    a_{n+1}/a_n &= \frac{(n+1-t)^{n+1}}{(n+1)^{n+1}}\times\frac{n^n}{(n-t)^n} = \left(\frac{(n+1-t)n}{(n+1)(n-t)}\right)^n\times\frac{n+1-t}{n+1} \\
                    &= \left(1 + \frac{t}{(n+1)(n-t)}\right)^n\times\frac{n+1-t}{n+1} = \sum_{k=0}^n \binom{n}{k}\left(\frac{t}{(n+1)(n-t)}\right)^k \times \frac{n+1-t}{n+1} \\
                    &\geq \left(1 + \frac{nt}{(n+1)(n-t)}\right) \times \frac{n+1-t}{n+1} = \frac{(n^2+n-t)(n+1-t)}{(n+1)^2(n-t)} \\
                    &= \frac{(n+1)^2(n-t) + t^2}{(n+1)^2(n-t)} > 1
                \end{align*}
            \end{proof}
        % \subsection{$\left(1-\frac{1}{n}\right)^n$は単調増加}
        % 	\label{(1-1/n)^nは単調増加}
        % 	\begin{proof}
        % 		\quad\par
        % 		$a_n \coloneqq \left(1-\frac{1}{n}\right)^n$とする。
        % 		$a_0 < a_1$は明らか。以下$n\geq 2$とする。
        % 		\begin{align*}
        % 			a_{n+1}/a_n &= \frac{\left(\frac{n}{n+1}\right)^{n+1}}{\left(\frac{n-1}{n}\right)^n} = \left(\frac{n^2}{n^2-1}\right)^n\times\frac{n}{n+1} = \left(1+\frac{1}{n^2-1}\right)^n\times\frac{n}{n+1} \\
        % 			&= \frac{n}{n+1}\sum_{k=0}^n \binom{n}{k} \left(\frac{1}{n^2-1}\right)^k \geq \frac{n}{n+1}\times\left(1+\frac{n}{n^2-1}\right) = \frac{n^3+n^2-n}{n^3+n^2-n-1} > 1
        % 		\end{align*}
        % 	\end{proof}
    \section{極限}
        \subsection{\texorpdfstring{$\lim_{n\rightarrow\infty}\left(1+1/n+\sum_{k=2}^\infty\frac{c_k}{n^k}\right)^n = e\quad(\sum_{k=2}^\infty c_k\;\;\mathrm{converges})$}{lim\_\{n→∞\} (1 + 1/n + Σ\_\{k=2\}\string^∞ c\_k/n\string^k)\string^n = e (Σ\_\{k=2→∞\}c\_k converges)}}
            \begin{proof}
                \quad\par
                \[1+1/n+\sum_{k=2}^\infty\frac{c_k}{n^k} = 1+\frac{1}{n}\left(1+\sum_{k=2}^\infty\frac{c_k}{n^{k-1}}\right)\]
                であるから
                $\epsilon>0$を任意にとると、十分大きい$n\in\mathbb{N}$に対して
                \[\left(1+\frac{1-\epsilon}{n}\right)^n \leq \left(1+\frac{1}{n}\left(1+\sum_{k=2}^\infty\frac{c_k}{n^{k-1}}\right)\right)^n \leq \left(1+\frac{1+\epsilon}{n}\right)^n\]
                両辺で$n\rightarrow\infty$として従う。
            \end{proof}
        \subsection{\texorpdfstring{$\lim_{n\rightarrow\infty}\left(1+1/n+o(1/n)\right)^n = e$}{lim\_\{n→∞\} (1+1/n+o(1/n))\string^n = e}}
            \begin{proof}
                \quad\par
                $o(1/n)$の部分を$f(n)$とおく。
                \[ \left(1+\frac{1}{n}+f(n)\right)^n = (1+1/n)^n + \sum_{i=1}^n \combi{n}{i}(1+1/n)^{n-i}f(n)^i \]
                $\lim_{n\to\infty}(1+1/n)^n = e$の証明より、上式において$(1+1/n)^{n-i}\leq 3$であることがわかっているので
                \begin{align*}
                    \sum_{i=1}^n \combi{n}{i}(1+1/n)^{n-i}f(n)^i &\leq 3\sum_{i=1}^n \combi{n}{i}f(n)^i = 3\sum_{i=1}^n \frac{1}{i!}\frac{n(n-1)\dotsm(n-i+1)}{n^i}n^if(n)^i \\
                    &\leq 3\sum_{i=1}^n \frac{1}{i!}(nf(n))^i \leq 3nf(n)\sum_{i=1}^n \frac{1}{i!}(nf(n))^{i-1}
                \end{align*}
                $n$を十分大きくすれば$nf(n)\leq 1$なので
                \[ \sum_{i=1}aa^n \combi{n}{i}(1+1/n)^{n-i}f(n)^i \leq 3nf(n)\sum_{i=1}^n \frac{1}{i!} \leq 6nf(n) \to 0 \as n\to\infty \]
            \end{proof}
    \section{漸化式}
        \subsection{Fibonacci数列}
            \subsubsection{つがいの増殖からの漸化式の導出}
                第$n$世代のつがいの数を$a_n$とし、そのうち成熟しているものを$m_n$、未成熟のものを$y_n$とすると、$a_{n+1} - a_n = m_n$であり、同時に$m_n = m_{n-1} + y_{n-1}$なので結局$a_{n+1} = a_n + a_{n-1}$となる。
        \subsection{\texorpdfstring{$a_{n+1}=ra_n + p(n),\;(r\neq1 ,\;p(n):\;d\text{次多項式})$}{a\_\{n+1\} = ra\_n + p(n) (r≠1, p(n): d次多項式)}の一般項は\texorpdfstring{$a_n = Ar^n + q(n),\;(A:\text{ 定数},\;q(n):\;d\text{次多項式})$}{a\_n = Ar\string^n + q(n) (A: 定数, q(n): d次多項式)}}
            \begin{proof}
                \quad\par
                $p(n) = b_dn^d + \dotsb + b_1n + b_0n$とする。
                $a_n = Ar^n + q(n)\;\left(q(n) = c_dn^d + \dotsb + c_1n + c_0\right)$とおいて漸化式に代入すると次式を得る。
                \[ Ar^{n+1} + q(n+1) = Ar^{n+1} + rq(n) + p(n) \quad \therefore\; q(n+1) = rq(n) + p(n) \]
                両辺の$k$次の項の係数を比較すると右辺では$rc_k + b_k$、左辺では$\sum_{i=k}^d\binom{i}{k}c_i$であり、$r\neq 1$なので$d$次の係数から順に見てゆき、$c_d,c_{d-1},\dotsc,c_1,c_0$の順に決定できる。
                最後に初期条件$a_0$(或いは$a_1,a_2,\dotsc,$\;値が分かっている項なら何でも良い)の値から$A$を決定できる。
            \end{proof}
        % \begin{comment}
        %     \subsection{連立線形項間漸化式の添字の有効範囲}
        %         \begin{shadebox}
        %             $N\in\naturalNumbers,\; a\neq b$とする。数列${x_i}$に関する漸化式
        %             \[
        %                 \begin{cases}
        %                     x_{i+1} + a x_i = c_i \\
        %                     x_{i+1} + b x_i = d_i
        %                 \end{cases}
        %                 (1\leq i \leq N\textcolor{red}{-1})
        %             \]
        %             の解の添字の有効範囲が1つ拡張されて$1\leq i\leq\textcolor{red}{N}$になる必要十分条件は$c_N - d_N = a d_{N-1} - b c_{N-1}$である。
        %         \end{shadebox}
        %         \begin{proof}
        %             \quad\par
        %             上の連立方程式を$x_{i+1}$を消去する方針で解くと次になる。
        %             \[ x_i = \frac{c_i-d_i}{a-b} \quad (1\leq i \leq N\textcolor{red}{-1}) \]
        %             一方で$x_{i+1}$を消去する方針で解くと次になる。
        %             \[ x_{i+1} = \frac{a d_i - b c_i}{a-b} \quad (1\leq i \leq N\textcolor{red}{-1}) \quad \therefore x_i = \frac{a d_{i-1} - b c_{i-1}}{a-b} \quad (2\leq i \leq \textcolor{red}{N}) \]
        %         \end{proof}
        % \end{comment}